\subsection{An eleven-arrow strong pair with a three-arrow exit}

\begin{lemma}[Deleting a retract-supported vertex]
\label{C:retract-supported-vertex}
Let $C$ contain distinct objects $a,b$ such that
$\End_C(b)=\{1_b,e\}$ with $e^2=e\ne1_b$.
Suppose $e$ fixes every arrow incident to $b$ and another object:
$ep=p$ for $p:x\to b$, $pe=p$ for $p:b\to x$, $x\ne b$.
If there are $u:b\to a$, $v:a\to b$ with $vu=e$, then deleting
$b$ preserves the homotopy type of the classifying space.
\end{lemma}
\begin{proof}
Write $C'$ for the full subcategory omitting $b$, and $J:C'\to C$
for its inclusion.  Define $G:C\to C'$ by $G(b)=a$, by fixing
all retained objects and arrows, and by
\[
 G(1_b)=1_a,\quad G(e)=uv,\quad
 G(p:x\to b)=up,\quad G(q:b\to x)=qv.
\]
The element $uv$ is idempotent.  A composite through $b$ is
preserved because $qvu p=qep=qp$.  A composite through a retained
object and beginning and ending at $b$ cannot equal $1_b$:
it is fixed by multiplication by $e$, whereas $1_b$ is not.
It therefore equals $e$, and its image is $uv$ as required.
The other functor equations follow from the stated support
identities.  Thus $GJ=\id_{C'}$.  The components $1_x$ at retained
objects and $v$ at $b$ give a natural transformation
$JG\Rightarrow\id_C$.  Hence $BJ$ and $BG$ are homotopy inverses.
\end{proof}

\begin{proposition}\label{C:eleven-strong-three-exit}
Let $S$ be the eleven-arrow two-object category formed by a
rectangular minimal ideal with three total rows and three total
columns, and one additional identity at each object.  Suppose
each object has at least one row and one column.  Form $C$ by
adjoining a sink $c$ with $\End(c)=\{1_c\}$ and three total
arrows from $S$ to $c$, with arbitrary compatible actions.
Then $BC$ does not have the homotopy type of $\mathbb CP^2$.
\end{proposition}
\begin{proof}
Write the two objects as $a,b$.  The row counts are $2,1$, and
the column counts are $2,1$.  The quotient of the rational
category algebra by its minimal ideal is
$\mathbb Q1_a\oplus\mathbb Q1_b$.  The natural categorical
balanced-tensor presentation therefore identifies $H_2(BS;\mathbb Q)$
with the tensor product of the row and column augmentation modules
over this semisimple algebra.  If the two-dimensional row and
column sets occur at different objects, the augmentation modules
have disjoint supports, so $H^2(BS;\mathbb Q)=0$.
Every degree-two class of the act cone is then relative to $BS$;
its products vanish because the relative quotient is the suspension
of the action category.  This excludes the indicated case.

Suppose instead that both two-dimensional sets occur at $a$.
Then $\End_S(a)=P_5$, the five-element sphere monoid, while
$\End_S(b)=\{1_b,e\}$ and both cross-hom sets have two elements.
The same balanced-tensor formula shows that the full corner
inclusion induces an isomorphism
\[
 H_2(BP_5;\mathbb Q)\longrightarrow H_2(BS;\mathbb Q).
                                                               \tag{\ref{C:eleven-strong-three-exit}.1}
\]
Both source objects have at least one arrow to $c$, since they
are in the same strongly connected component.  Their arrow
counts are therefore either $(2,1)$ or $(1,2)$.

In the first case, the unique arrow $b\to c$ is fixed by $e$,
as are all the original arrows between $a$ and $b$.
Any $u:b\to a$, $v:a\to b$ satisfy $vu=e$ by the rectangular
composition law.  Lemma~\ref{C:retract-supported-vertex} deletes
$b$.  The resulting category is an act cone over $P_5$ with two
states and has eight arrows.  Its degree-two products vanish
by Theorem~\ref{C:two-row-act-cups}.

In the second case there is a unique arrow $a\to c$, fixed by
every element of $P_5$.  It gives a functor from $P_5$ to the
action category $E$ of the exit, whose composite with
$E\to S$ is the full-corner inclusion.  By
(\ref{C:eleven-strong-three-exit}.1), the map
$H_2(BE;\mathbb Q)\to H_2(BS;\mathbb Q)$ is surjective.
Thus $H^2(BS;\mathbb Q)\to H^2(BE;\mathbb Q)$ is injective.
The cofiber exact sequence again makes every degree-two class
of $BC$ relative to $BS$, and the relative-suspension argument
makes every degree-two product zero.  All cases contradict the
nonzero square of the rational generator on $\mathbb CP^2$.
\end{proof}
